% $Header: /cvsroot/latex-beamer/latex-beamer/solutions/generic-talks/generic-ornate-15min-45min.en.tex,v 1.4 2004/10/07 20:53:08 tantau Exp $

\documentclass{beamer}

% This file is a solution template for:

% - Giving a talk on some subject.
% - The talk is between 15min and 45min long.
% - Style is ornate.



% Copyright 2004 by Till Tantau <tantau@users.sourceforge.net>.
%
% In principle, this file can be redistributed and/or modified under
% the terms of the GNU Public License, version 2.
%
% However, this file is supposed to be a template to be modified
% for your own needs. For this reason, if you use this file as a
% template and not specifically distribute it as part of a another
% package/program, I grant the extra permission to freely copy and
% modify this file as you see fit and even to delete this copyright
% notice. 


\mode<presentation>
{
  \usetheme{Antibes}
  \useinnertheme{circles}
  \setbeamercovered{transparent}
  % or whatever (possibly just delete it)
}


\usepackage[english]{babel}
% or whatever

% \usepackage[latin1]{inputenc}
\usepackage[utf8]{inputenc}
% or whatever

\usepackage{times}
\usepackage[T1]{fontenc}
% Or whatever. Note that the encoding and the font should match. If T1
% does not look nice, try deleting the line with the fontenc.

\usepackage{multimedia}

\usepackage{graphicx}

\usepackage{colortbl}

\usepackage[sectionbib, round, authoryear, longnamesfirst]{natbib}

\title[Improvisation, Computers, and Primary Process]% (optional, use only with long paper titles)
{Improvisation, Computers, and Primary Process: Why improvise with computers?}

\author % (optional, use only with lots of
% authors)
{Henrik Frisk}
%{Henrik Frisk \and Stefan
%\"{O}stersj\"{o}\\ \texttt{henrik.frisk@mhm.lu.se}, \texttt{stefan_ostersjo@hotmail.com}}
% - Use the \inst{?} command only if the authors have different
%   affiliation.

\institute[MHM] % (optional, but mostly needed)
{
  Malm\"{o} Academy of Music, Lund University\\
  http://www.mhm.lu.se\\
  http://www.henrikfrisk.com
}
% - Use the \inst command only if there are several affiliations.
% - Keep it simple, no one is interested in your street address.

\date[Artistic research...] % (optional)
{\tiny{ARTISTIC RESEARCH IN CONTEMPORARY MUSIC: COMPOSITION – INTERPRETATION - IMPROVISATION}\\[0.1cm]Odense, Oct. 29, 2008}

\subject{Talks}
% This is only inserted into the PDF information catalog. Can be left
% out. 

% If you have a file called "university-logo-filename.xxx", where xxx
% is a graphic format that can be processed by latex or pdflatex,
% resp., then you can add a logo as follows:

% \pgfdeclareimage[height=1cm]{university-logo}{../img/mhm-logo}
% \logo{\pgfuseimage{university-logo}}

% \beamerdefaultoverlayspecification{<+->}

\begin{document}

\begin{frame}
  \titlepage
  \includegraphics[height=1.5cm]{../img/mhm-logo}
\end{frame}

\section{Introduction}

\begin{frame}
  \frametitle{Artistic Research}
    \begin{figure}[h]
    \centering
    \includegraphics<2>[width=10em]{../img/forskning}
  \end{figure}
\end{frame}

  \begin{frame}
    \frametitle{Improvisation,  Computers and Interaction}
      \begin{figure}[h]
    \centering
    \includegraphics[width=12em]{../img/ImpCompInt-cover}
  \end{figure}
  \end{frame}

\section{Improvisation: definition?}

\begin{frame}
  \frametitle{What is Improvisation?}
  % - A title should summarize the slide in an understandable fashion
  %   for anyone how does not follow everything on the slide itself.

  \begin{block}
    
  \begin{itemize}

  \item<1-> A substitute for organization?

  \item<2- | alert@4> Bruno Nettl: ``[i]n the conception of the art music world,
    improvisation embodies the absence of precise planning and
    discipline'' {\footnotesize(\emph{An art neglected in
        scholarship})}

  \item<3-> Bruno Nettl: ``discipline, intricacy, and control of
    complexities all play major roles in various kinds of
    improvisation'' {\footnotesize (\emph{Ibid.})}

  \item<5-> But, improvisation is \emph{also} absence of control.

  \end{itemize}
  \end{block}
\end{frame}

\begin{frame}
  \frametitle{Basis for current discussion}
  \begin{block}{The nature of improvisation}
    \begin{itemize}
    \item<1-> Lack of control may in fact be important features.
    \item<2-> Computer technology may help to `un-learn' stylistically bound musical knowledge.
    \item<3-> The `un-learning' may be understood in terms of the concept of \emph{primary process}.
    \end{itemize}
  \end{block}
\end{frame}

\section{Computers and Improvisation}

\begin{frame}
  \frametitle{Computers and Improvisation}
  \begin{block}<1->{Computers-as-improvisers}
    \begin{itemize}
    \item<2-> Computers are poor improvisers.
    \item<3-> The human capacity for making ``good choices'' spontaneously is enormous
      and still close to impossible to model in a computer program.
    \end{itemize}
  \end{block}
  \begin{block}<4->{Computers-as-instruments}
    \begin{itemize}
    \item<5-> A human `computer-player' is the interface.
    \item<6-> Any sound is possible: dream or disaster?
    \end{itemize}
  \end{block}
\end{frame}

\begin{frame}
  \frametitle{Computer and Improvisation}
  \begin{block}{Summary}
    \begin{itemize}
    \item<1-> Non-intuitive musical UI.
    \item<2-> No physical contact between body and sound.
    \item<3-> No unified notion of what the computer-as-instrument consist of.
    \item<4-> The computer is useless as an instrument for improvisation(?)
    \item<5-| alert@5> Yet, the computer is an agent in the practice of improvisation.
    \end{itemize}
  \end{block}
\end{frame}

\section{Improvisation and the Mind}


\begin{frame}
  \frametitle{Improvisation and sensibility}
  \begin{block}{Personhood}
    \begin{itemize}
    \item<1-> G. Lewis: ``Afrological sensibility''
    \item<2-> G. Russel: ``Intuitive Intelligence''
    \item<3-> The personal sound.
    \item<4-> Conclusion: Improvised music thrives on sensibility.
    \item<5-> In what sense may the computer be an agent of sensibility?
    \end{itemize}
  \end{block}
\end{frame}

\begin{frame}
  \frametitle{Primary process}
  \begin{block}{The language of the unconscious}
    \begin{itemize}
    \item<1-> Bateson: The Algorithms of the heart.
    \item<2-> Music as a method of translation between primary and secondary process.
    \item<3-> Habit: actions become unconscious.
    \item<4-> Habit formation is not a guarantee for access to the unconscious.
    \end{itemize}
  \end{block}
\end{frame}

\section{In practice}


\begin{frame}
  \frametitle{Ornette Coleman}
  \begin{block}{Picking up the violin and the trumpet.}
    \begin{itemize}
    \item<1-> To rid himself of the Parker legacy.
    \item<2-> Striving to play ``without memory''
    \item<3-> Free from extra-musical parameters: Denardo Coleman
    \item<4-> Unknown instruments gave him internal freedom, liberated from physical memory.
    \end{itemize}
  \end{block}
\end{frame}

\begin{frame}
  \frametitle{Marcel Duchamp}
  \begin{block}{\emph{The Bride Stripped Bare by Her Bachelors, Even
        (The Large Glass)} (1915-1923)}
    \begin{itemize}
    \item<1-> To rid himself of old fashioned drawing.
    \item<2-> Knowledge reversed but not forgotten: Memory as a function of the future''
    \item<3-> Duchamp: ``I unlearned to draw. The point was to forget
      \emph{with my hand}.''
    \item<4-> To consciously unlearn: Is it at all possible?
    \end{itemize}
  \end{block}
\end{frame}

\section{Why computers?}

\begin{frame}
  \frametitle{Resistance}
  \begin{block}{Why forget?}
    \begin{itemize}
    \item<1-> To \emph{increase} resistance.
    \item<2-> To oppose habit.
    \item<3-> The computer is an agent in improvised music not despite its shortcomings but \emph{thanks to them}
    \item<4-> Frustration of an instrument's shortcomings are secondary process. 
    \item<5-> Playing it is primary process.
    \end{itemize}
  \end{block}
\end{frame}

\section{Conclusion}

\begin{frame}
  \frametitle{Summary}
  \begin{block}<1->{The computer as a forgetting machine}
    \begin{itemize}
    \item<2-> Musicians and artist have developed methods to achieve
      that which the computer allows us do prior to its existence.
    \item<3-> Habit formation is economy of conscious thought rather than access to unconscious.
    \end{itemize}
  \end{block}
  \begin{block}<4->{Aden Evens and the body}
    \begin{itemize}
    \item<5-> ``The musician who just allows his body to do what it
      will cannot be immerse in the music [\ldots] the one who lets his body play is not putting
      himself on the line''
    \end{itemize}
  \end{block}
\end{frame}

\begin{frame}
  \frametitle{Thank you}
I wish to thank the Center for Ny Musik for the invitation.

\end{frame}

\begin{frame}
  \frametitle{Questions?}
\includegraphics<1>[width=13em]{../img/QuestionMark}  
\end{frame}
\end{document}


